\chapter{PENDAHULUAN}

%Latar belakang 
\section{Latar Belakang}
Pelaporan masalah lingkungan oleh masyarakat saat ini masih menjadi kesulitan karena cara yang digunakan oleh pemerintah daerah kebanyakan masih menggunakan \textit{Short Message Service} (SMS) atau melalui whatsapp. Meskipun sudah ada yang menggunakan website namun, berdasarkan analisis data yang dikumpulkan oleh tim peneliti, untuk merespon satu laporan diperlukan rata-rata waktu hampir dua jam hanya untuk menugaskan laporan yang masuk kepada instansi terkait . Hal tersebut terjadi karena pemerintah daerah yang harus secara manual untuk menentukan instansi yang tepat untuk menangani laporan masyarakat yang masuk, seperti sampah menumpuk, banjir, pohon tumbang, kerusakan fasilitas umum, dan pencemaran. Masalah lingkungan seperti ini bisa cepat diatasi jika masyarakat langsung yang melaporkan, karena masyarakat langsung yang ada di tempat kejadian. Jika masyarakat melaporkan melalui website dengan sistem yang bagus maka Dinas Lingkungan atau Instansi yang sesuai dengan kejadian bisa langsung datang ke lapangan dan mengatasi masalah sehingga aktivitas masyarakat dapat berjalan lancar kembali. Jika sistem pelaporan yang hanya menggunakan teks atau formulir, sistem pelaporan tersebut masih memiliki kelemahan seperti kurangnya informasi visual yang dapat memastikan bahwa laporan tersebut valid karena hal seperti ini yang dapat menyebabkan keterlambatan penanganan atau kesalahan dalam prioritas penanganan masalah\citep{siretPublicComplainingBlessing2022}  

Integrasi data multimodal berupa gambar dan teks dalam sistem pelaporan menawarkan solusi yang efektif untuk mengatasi permasalahan tersebut. Pendekatan multimodal memungkinkan sistem untuk memanfaatkan kekuatan dari kedua modalitas: gambar memberikan bukti visual kondisi lapangan yang objektif, sementara teks menyediakan informasi kontekstual seperti deskripsi detail, lokasi, dan tingkat kerusakan. 
Dengan mengimplementasikan \textit{Artificial Intelligence} (AI) berbasis arsitektur transformer, data gambar dapat diklasifikasikan secara otomatis sesuai dengan jenis permasalahan lingkungan, sedangkan data teks diproses untuk mengekstraksi informasi penting yang mendukung pengambilan keputusan \citep{baltrusaitisMultimodalMachineLearning2019}. Kombinasi kedua modalitas ini tidak hanya meningkatkan akurasi klasifikasi, tetapi juga memperkaya konteks informasi yang diperlukan untuk penugasan laporan kepada instansi yang tepat secara otomatis dan efisien \citep{oquabDINOv2LearningRobust2024}. Pendekatan multimodal ini diharapkan dapat mengurangi waktu respons secara signifikan dan meningkatkan kualitas penanganan masalah lingkungan secara keseluruhan. 

Sistem pelaporan ini juga dapat membantu proses pengambilan keputusan yang lebih tepat sasaran. Maka pembuatan sistem ini akan mendukung terbentuknya ekosistem smart city yang lebih efisien, partisipatif, dan adaptif terhadap kebutuhan masyarakat, maka penelitian mengenai ”PENGEMBANGAN SISTEM PELAPORAN MASALAH LINGKUNGAN DENGAN INTEGRASI DATA GAMBAR DAN TEKS UNTUK MENDUKUNG SMART CITY” menjadi penting untuk dilakukan agar menghasilkan solusi yang inovatif dalam manajemen masalah lingkungan yang lebih cerdas, cepat, dan responsif. 

%rumusan masalah
\section{Rumusan Masalah}
Berdasarkan latar belakang di atas, penelitian ini akan menjawab beberapa pertanyaan utama: 
\begin{enumerate}
    \item Bagaimana merancang dan membangun sistem pelaporan masalah lingkungan berbasis multimodal yang mengintegrasikan data gambar dan teks deskripsi secara efektif? 
    \item Bagaimana mengimplementasikan model klasifikasi berbasis arsitektur transformer untuk mengkategorikan jenis masalah yang dilaporkan masyarakat dari kombinasi data gambar dan teks deskripsi?
    \item Sejauh mana integrasi gambar dan teks dapat meningkatkan kualitas informasi laporan dibandingkan teks saja?  
    \item Bagaimana merancang mekanisme penugasan \textit{(assignment)} otomatis berbasis \textit {artificial intelligence} yang dapat mendistribusikan laporan kepada instansi terkait secara akurat dan efisien? 
\end{enumerate}

%Ruang Lingkup 
\section{Ruang Lingkup}
Ruang lingkup penelitian ini adalah sebagai berikut:
\begin{enumerate}
    \item Jenis masalah 
        \begin{enumerate}
            \item Batasan penelitian ini berfokus pada degradasi lingkungan yang terlihat seperti  jalanan rusak, kecelakaan, kerusakan trotoar, vandalisme, sampah, kerusakan halte, kecelakaan lalu lintas, dan pohon tumbang.
            \item Tidak mencakup masalah lingkungan yang membutuhkan pengukuran khusus seperti polusi udara dan pencemaran air secara kimia. 
        \end{enumerate}
    \item Data 
        \begin{enumerate}
            \item Data yang digunakan berbentuk gambar kondisi lingkungan dan teks berupa lokasi masalah lingkungan, deskripsi singkat laporan dari masyarakat, dan keterangan tambahan. 
            \item Data diambil dari dataset yang ada di internet atau hasil pengumpulan secara manual melalui \textit{web scraping}. 
        \end{enumerate}
    \item Fokus Penelitian 
        \begin{enumerate}
            \item Sistem hanya difokuskan untuk proses pelaporan, penyimpanan, dan klasifikasi masalah lingkungan berdasarkan gambar dan teks. 
        \end{enumerate}
    \item Fitur utama
        \begin{enumerate}
            \item User dapat melakukan upload gambar/foto masalah lingkungan, mengisi teks deskripsi. 
            \item Menampilkan dasbor peta interaktif untuk admin mengetahui dengan mudah lokasi masalah lingkungan. 
            \item Analisis gambar menggunakan AI untuk mengidentifikasi kategori masalah dan analisis teks untuk memperjelas kondisi masalah 
        \end{enumerate}
    \item user 
        \begin{enumerate}
            \item Masyarakat sebagai pelapor masalah lingkungan 
            \item Admin dari instansi terkait yang bertugas menindaklanjuti laporan masyarakat. 
        \end{enumerate}
    \ Limitasi Teknologi 
        \begin{enumerate}
            \item Jenis masalah lingkungan yang dapat ditangani terbatas, hanya yang mudah terlihat oleh mata saja, tidak mencakup masalah yang harus menggunakan alat khusus seperti polusi udara atau pencemaran air. 
            \item Sistem yang dibangun hanya berbentuk prototipe sehingga tidak mencakup integrasi langsung dengan pembangunan smart city atau penanganan lanjut masalah di lapangan. 
        \end{enumerate}
 \end{enumerate}

%Tujuan dan Manfaat Penelitian 
 \section{Tujuan dan Manfaat Penelitian}
 
 \begin{enumerate}
    \item Tujuan Penelitian 
    
    Berdasarkan rumusan masalah yang diuraikan, penelitian ini memiliki beberapa tujuan utama. Berikut adalah beberapa tujuan yang ingin dicapai dalam penelitian ini: 
        \begin{enumerate}
            \item Mengembangkan sistem pelaporan masalah lingkungan berbasis web yang mengintegrasikan data multimodal (gambar dan teks). 
            \item Mengimplementasikan klasifikasi otomatis menggunakan arsitektur \textit {transformer} untuk mengkategorikan masalah dan menugaskan laporan kepada instansi terkait.
            \item Merancang dan memnguban dasbor visualisasi data untuk memfasilitasi pengambilan keputusan dan pemantauan secara \textit{real-time}.  
        \end{enumerate}

    \item Manfaat Penelitian
    \begin{enumerate}
        \item Meningkatkan Pelayanan Publik: Sistem ini diharapkan dapat mempercepat proses pelaporan dan penanganan masalah, sehingga meningkatkan kualitas pelayanan publik.
        \item Integrasi Data Multimodal: Penelitian ini memberikan manfaat pada pengembangan sistem klasifikasi yang lebih robust melalui penggunaan model \textit{state-of-the-art} (DINOv2 \& Multilingual E5). Manfaat utamanya adalah peningkatan akurasi identifikasi laporan melalui sinergi data teks dan gambar, yang meminimalisir kesalahan klasifikasi
    \end{enumerate}

 \end{enumerate}

 \section{Metode Penelitian}
    Metode penelitian yang digunakan dalam penelitian ini terdiri dari beberapa tahap utama, yaitu:
    
    \subsection{Pengumpulan data}
     Data yang akan digunakan dalam penelitian ini diperoleh melalui dua metode, yaitu pengambilan data dari dataset publik yang tersedia di platform Kaggle dan pengumpulan data secara mandiri melalui metode \textit{web scraping} pada portal pengaduan masyarakat crm.jakarta.go.id. Data yang dikumpulkan meliputi gambar kondisi lingkungan dan teks deskripsi laporan dari masyarakat.
    \subsection{Pengembangan  \textit{Artificial Intelegence}}
    Tahap ini merupakan inti dari eksperimen AI, di mana pengembangan model difokuskan pada penggabungan fitur dari data gambar dan teks.Selanjutnya dilakukan pemilihan dan perancangan model melalui eksplorasi dua pendekatan utama yaitu menggunakan \textit{Transformer} dan konvensional. 
    \subsection {Pengembangan  \textit{Website}}
    Metode implementasi aplikasi web ini berfokus pada prosedur deployment model yang telah terlatih ke dalam lingkungan produksi dengan membungkusnya dalam API. Melalui antarmuka pengguna yang responsif, sistem memungkinkan pengunggahan data laporan, pemrosesan di sisi server, hingga penayangan hasil prediksi dan dasbor visualisasi secara instan. Integrasi ini bertujuan untuk mentransformasi model AI menjadi alat praktis yang dapat diakses melalui \textit{browser} guna mendukung pemantauan laporan masyarakat. 